\documentclass[12pt,a4paper]{article}

% =====================================================
% PACOTES
% =====================================================
\usepackage[utf8]{inputenc}
\usepackage[T1]{fontenc}
\usepackage[brazil]{babel}
\usepackage{geometry}
\usepackage{graphicx}
\usepackage{fancyhdr}
\usepackage{titlesec}
\usepackage{xcolor}
\usepackage{hyperref}
\usepackage{setspace}
\usepackage{enumitem}
\usepackage{newtxtext,newtxmath}
\usepackage{etoolbox}
\usepackage{multicol}
\usepackage{tabularx}

% Toggle para o Gabarito (false = prova aluno, true = gabarito professor)
\newtoggle{gabarito}
\settoggle{gabarito}{true}

\geometry{
    top=1.5cm,
    bottom=2cm,
    left=1.5cm,
    right=1.5cm
}

\onehalfspacing

% =====================================================
% CORES PERSONALIZADAS
% =====================================================
\definecolor{azulTitulo}{RGB}{30,60,90}
\definecolor{corGabarito}{RGB}{180,50,50}

% --- COMANDOS PARA O GABARITO ---
\newcommand{\correto}[1]{\iftoggle{gabarito}{\textcolor{corGabarito}{\textbf{#1}}}{#1}}
\newcommand{\justificativa}[1]{%
\vspace{0.1cm}\par\noindent\framebox[\linewidth]{%
\parbox[c][1.2cm][t]{\dimexpr\linewidth-2\fboxsep-2\fboxrule\relax}{%
\vspace{0.1cm}\textbf{Justificativa para falsa: }\iftoggle{gabarito}{\textcolor{corGabarito}{#1}}{}\vspace{0.1cm}%
}%
}%
}
\newcommand{\resp}[1]{\iftoggle{gabarito}{\textcolor{corGabarito}{\textbf{#1}}}{\phantom{#1}}}

\begin{document}

\noindent
\makebox[\textwidth]{Aluno(a): \dotfill Data: \_\_/\_\_/\_\_\_\_}

\vspace{0.3cm}
\begin{center}
    \textbf{\Large Prova de Recuperação Teórica de Programação para Internet 1}\\
    \textbf{CST em Sistemas para Internet}
\end{center}

\noindent
\textbf{Orientações:} Prova individual. Responda com clareza. Nas questões de Verdadeiro ou Falso, justifique as falsas.

\vspace{0.3cm}
\noindent
\textbf{Gabarito.} Preencha cada célula com a letra da opção correta, ou V/F/Sequência.

\vspace{0.2cm}
\begin{center}
\renewcommand{\arraystretch}{1.5}
\begin{tabular}{|*{10}{c|}}
\hline
01 & 02 & 03 & 04 & 05 & 06 & 07 & 08 & \hspace{0.4cm}09\hspace{0.4cm} & 10 \\
\hline
\iftoggle{gabarito}{
 \textcolor{corGabarito}{B} & \textcolor{corGabarito}{A} & \textcolor{corGabarito}{C} & \textcolor{corGabarito}{C} & \textcolor{corGabarito}{F} & \textcolor{corGabarito}{V} & \textcolor{corGabarito}{F} & \textcolor{corGabarito}{V} & \textcolor{corGabarito}{2143} & \textcolor{corGabarito}{B}
}{
 & & & & & & & & & 
} \\
\hline
\end{tabular}
\end{center}
\vspace{0.1cm}
\noindent\small\textit{* Questão 9 é de relacionar colunas.}

\vspace{0.4cm}

\begin{multicols}{2}

\textbf{1.} Qual tag semântica do HTML5 deve ser utilizada para agrupar links de navegação interna ou externa, definindo a estrutura de menus ou sumários do site?
\begin{enumerate}[label=(\alph*)]
    \item \texttt{<section>}
    \item \correto{\texttt{<nav>}}
    \item \texttt{<header>}
    \item \texttt{<aside>}
\end{enumerate}

\vspace{0.2cm}
\textbf{2.} Ao trabalhar com o modelo Flexbox no CSS, qual propriedade deve ser aplicada ao elemento pai (container) para fazer com que os itens filhos quebrem para a linha seguinte caso não caibam no espaço horizontal disponível?
\begin{enumerate}[label=(\alph*)]
    \item \correto{\texttt{flex-wrap: wrap;}}
    \item \texttt{flex-direction: column;}
    \item \texttt{justify-content: space-between;}
    \item \texttt{display: grid;}
\end{enumerate}

\vspace{0.2cm}
\textbf{3.} No JavaScript, ao selecionarmos o elemento de entrada de texto e salvá-lo na variável \texttt{campo}, qual propriedade acessa o valor digitado?
\begin{enumerate}[label=(\alph*)]
    \item \texttt{campo.textContent}
    \item \texttt{campo.innerHTML}
    \item \correto{\texttt{campo.value}}
    \item \texttt{campo.type}
\end{enumerate}

\vspace{0.2cm}
\textbf{4.} Para que serve a função global \texttt{fetch()} no JavaScript moderno?
\begin{enumerate}[label=(\alph*)]
    \item Para formatar e aplicar estilos CSS de forma automática na página.
    \item Para declarar variáveis locais constantes.
    \item \correto{Para realizar requisições HTTP assíncronas sem recarregar a página.}
    \item Para criar um banco de dados local na memória do navegador.
\end{enumerate}

\vspace{0.2cm}
\textbf{5. [V/F]} ( ) A propriedade de medida relativa \texttt{em} no CSS calcula seu valor com base no tamanho da fonte do elemento raiz do documento (a tag \texttt{<html>}), sendo ideal para criar dimensionamentos globais responsivos.
\justificativa{Falso. A unidade relativa ao elemento raiz é o \texttt{rem} (\textit{root em}). O \texttt{em} é relativo ao tamanho de fonte do próprio elemento ou do elemento pai imediato.}

\vspace{0.2cm}
\textbf{6. [V/F]} ( ) No JavaScript, podemos alterar a cor de fundo de um elemento no DOM modificando a propriedade \texttt{style.backgroundColor} (ex: \texttt{el.style.backgroundColor = "blue"}).
\justificativa{Verdadeiro.}

\vspace{0.2cm}
\textbf{7. [V/F]} ( ) A função \texttt{fetch()} no JavaScript rejeita automaticamente a sua Promise correspondente (caindo no bloco \texttt{.catch()}) sempre que o servidor responde com um código de erro HTTP do tipo \texttt{404} (Não Encontrado) ou \texttt{500} (Erro Interno).
\justificativa{Falso. O \texttt{fetch()} só rejeita a Promise em caso de falha de rede/conexão ou bloqueio de CORS. Erros HTTP como 404 e 500 resolvem a Promise com sucesso, sendo necessário verificar a propriedade \texttt{response.ok} ou \texttt{response.status} no \texttt{.then()}.}

\vspace{0.2cm}
\textbf{8. [V/F]} ( ) O evento do tipo \texttt{submit} no JavaScript é disparado em um elemento de formulário (\texttt{<form>}) sempre que o usuário tenta enviá-lo, por exemplo, ao clicar em um botão \texttt{<button type="submit">} contido no formulário.
\justificativa{Verdadeiro.}

\vspace{0.2cm}
\textbf{9. [Sequência]} Associe a Coluna A (Seletor/Pseudo-classe CSS) com a Coluna B (Descrição de funcionamento).
\begin{enumerate}[label=(\arabic*), leftmargin=*]
    \item \texttt{.card:hover}
    \item \texttt{div > p}
    \item \texttt{input[type="submit"]}
    \item \texttt{\#destaque}
\end{enumerate}
\vspace{-0.2cm}
\textbf{Coluna B:}
\begin{itemize}[leftmargin=*, label={}]
    \item ( \resp{2} ) Seleciona qualquer parágrafo (\texttt{<p>}) que seja um filho direto de um container \texttt{<div>}.
    \item ( \resp{1} ) Aplica estilos a um elemento da classe \texttt{card} somente quando o cursor do mouse estiver posicionado sobre ele.
    \item ( \resp{4} ) Seleciona o único elemento da página que possui a identificação exclusiva \texttt{id="destaque"}.
    \item ( \resp{3} ) Seleciona elementos de entrada do tipo formulário cujo atributo \texttt{type} seja exatamente "submit".
\end{itemize}

\end{multicols}

\vspace{0.3cm}
\hrule
\vspace{0.3cm}
\textbf{10.} Analise o trecho de código JavaScript abaixo:

\begin{verbatim}
const botaoAdd = document.querySelector('#btn-adicionar');
const inputTarefa = document.querySelector('#nova-tarefa');
const listaTarefas = document.querySelector('#lista');

botaoAdd.addEventListener('click', () => {
    const texto = inputTarefa.value.trim();
    
    if (texto !== '') {
        const item = document.createElement('li');
        item.textContent = texto;
        
        item.addEventListener('click', () => {
            item.classList.toggle('concluida');
        });
        
        listaTarefas.appendChild(item);
        inputTarefa.value = '';
    }
});
\end{verbatim}

Supondo que no HTML existam os elementos com seletores corretos, qual será o comportamento do sistema quando um usuário digitar "Estudar JS", clicar no botão de adicionar e, em seguida, clicar sobre o texto adicionado na lista?
\begin{enumerate}[label=(\alph*)]
    \item O texto "Estudar JS" será adicionado à lista e, ao ser clicado, será removido de forma definitiva do DOM da página.
    \item \correto{O texto "Estudar JS" será inserido na lista e, ao receber o clique, terá a classe CSS "concluida" alternada (adicionada ou removida), possibilitando estilização visual dinâmica.}
    \item O código apresentará um erro de execução porque o evento de clique no item de lista (\texttt{li}) foi registrado após a criação do elemento.
    \item O texto será adicionado, mas o campo de input de texto não será limpo.
\end{enumerate}

\vfill
\begin{center}
    \footnotesize \textit{IFSC Campus Garopaba -- Prof. André Moraes}
\end{center}

\end{document}
